\documentclass[11pt]{amsart}

\usepackage[T1]{fontenc}
\usepackage{lmodern}
\usepackage{microtype}
\usepackage{mathtools,amssymb,amsthm}
\usepackage[hidelinks]{hyperref}

\hypersetup{
  pdftitle={Four Mutually Unbiased Bases in Dimension Eight Give a
Non-Decomposable Witness}
}

\newtheorem{theorem}{Theorem}
\newtheorem{lemma}[theorem]{Lemma}
\newtheorem{corollary}[theorem]{Corollary}

\DeclareMathOperator{\Tr}{tr}
\newcommand{\CC}{\mathbb C}
\newcommand{\FF}{\mathbb F_2}
\newcommand{\W}{W^{\Gamma}}
\newcommand{\M}{\mathcal M}
\newcommand{\ii}{\mathrm i}

\title[A non-decomposable four-MUB witness in dimension eight]
{Four Mutually Unbiased Bases in Dimension Eight Give a Non-Decomposable
Witness}

\date{}

\subjclass[2020]{81P40, 81P55, 15B48}
\keywords{mutually unbiased bases, bound entanglement, entanglement witness,
decomposable witness, positive partial transpose, counterexample}

\begin{document}

\begin{abstract}
For $d=2^r$ with $r>1$, Bae, Bera, Chru\'sci\'nski, Hiesmayr and McNulty conjectured
that the minimal number of mutually unbiased bases needed to detect bound entanglement
is $m=\tfrac d2+1$. We refute this at $d=8$, inside the one witness family that
statement quantifies over. We exhibit a set $\M_4$ of four mutually
unbiased bases of $\CC^8$ --- so $m=\tfrac d2$, one below the conjectured minimum ---
and a shift $s=1$, for which the witness $\W(\M_4,s)$ of their family is
\emph{non-decomposable}. The certificate is a positive partial transpose state $\rho$
on $\CC^8\otimes\CC^8$, printed here in full as an integer matrix divided by
$1049602$, with
\[
  \Tr\bigl[\W(\M_4,1)\,\rho\bigr]=-\frac{14588}{524801}<0 .
\]
Every quantity is exact: $64\,\W(\M_4,1)$ is an integral symmetric matrix, and the
positivity of $\rho$ and of $\rho^{\Gamma}$ is certified by exact rational
$LDL^{\!\top}$ factorisations. Hence the least number of mutually unbiased bases of
$\CC^8$ for which that family has a non-decomposable member is at most $4$, and not
the $\tfrac d2+1=5$ asserted at $r=3$.
\end{abstract}

\maketitle

\section{The conjecture}

Let $\M_m=\{B^{(1)},\dots,B^{(m)}\}$ be a set of $m$ pairwise mutually unbiased bases
(MUBs) of $\CC^d$. Spengler, Huber, Brierley, Adaktylos and Hiesmayr
\cite{Spengler} attach to $\M_m$ and a shift $s\in\mathbb Z_d$ an operator
$\W(\M_m,s)$ on $\CC^d\otimes\CC^d$ (recalled in \S2) which is nonnegative on every
separable state. When $\W(\M_m,s)$ is \emph{non-decomposable}, i.e.\ cannot be written
as $P+Q^{\Gamma}$ with $P,Q\succeq0$, it detects a state with positive partial
transpose, that is a bound entangled state. Bae, Bera, Chru\'sci\'nski, Hiesmayr and
McNulty \cite{BBCHM} asked how many MUBs are needed for this, and conjectured a
sharp answer in the power-of-two dimensions.

Their Conjecture~1 reads: ``\emph{If $d=2^r$ and $r>1$, the minimal number of MUBs
required for bound entanglement detection equals $m=\frac d2+1=2^{r-1}+1$}.''

Two clauses of \cite{BBCHM} fix what a counterexample must be. The sentence
immediately above the conjecture reads ``\emph{When $d=2^r$, we predict that
$\frac d2+1$ MUBs are both necessary and sufficient for the witness
$\mathbf W^{\Gamma}(\mathcal M_m,s)$ to be non-decomposable}'', so the minimum is taken
over the paper's own family and ``detects bound entanglement'' means ``is
non-decomposable''. Second, the minimum is \emph{existential}: at $d=4$ that paper
exhibits particular triples of MUBs and still calls $\frac d2+1$ minimal. One
$m=\frac d2$ example at $d=8$ is therefore the right shape of counterexample, and
everything below stays inside that family.

What is settled below is the cell $r=3$, by one exhibited object: at $d=8$ the least
$m$ for which the family \eqref{eq:W} has a non-decomposable member is at most
$4=\frac d2$, so it is not $\frac d2+1=5$. Only the necessity half of the quoted
statement is at issue; the sufficiency theorem of \cite{BBCHM} concerns
$m>\frac d2+1$.

\section{The witness family}

Fix $d=8$ and index the computational basis of $\CC^8$ by
$x=4x_1+2x_2+x_3$ with $x_1,x_2,x_3\in\FF$. For a symmetric matrix
$S\in\FF^{3\times3}$ and $v\in\FF^3$ put
\begin{equation}\label{eq:vS}
  |v_S\rangle
  =\frac1{\sqrt8}\sum_{x\in\FF^3}
   \ii^{\,S_{11}x_1+S_{22}x_2+S_{33}x_3}
   (-1)^{S_{12}x_1x_2+S_{13}x_1x_3+S_{23}x_2x_3+v\cdot x}\,|x\rangle ,
\end{equation}
where the exponent of $\ii$ is read modulo $4$ and the exponent of $-1$ modulo $2$,
and set $B_S=\{\,|v_S\rangle : v\in\FF^3\,\}$. Each $B_S$ is an orthonormal basis of
$\CC^8$. We encode $S$ by the integer
\begin{equation}\label{eq:idx}
  \mathrm{idx}(S)=S_{11}+2S_{12}+4S_{13}+8S_{22}+16S_{23}+32S_{33}\in\{0,\dots,63\}.
\end{equation}

Let $\M_m=\{B^{(1)},\dots,B^{(m)}\}$ with $B^{(1)}$ the computational basis, and let
$s\in\mathbb Z_8$. Writing $\overline{\,\cdot\,}$ for entrywise complex conjugation in
the computational basis, the witness of \cite{Spengler} in the form quantified over by
the conjecture is
\begin{equation}\label{eq:W}
\begin{split}
  \W(\M_m,s)
  =\ &\frac{d+m-1}{d}\,I\otimes I
   -\sum_{l=0}^{d-1}|l\rangle\langle l|\otimes|l\oplus s\rangle\langle l\oplus s|\\
   &-\sum_{k=2}^{m}\ \sum_{v}
     |v^{(k)}\rangle\langle v^{(k)}|\otimes
     \overline{|v^{(k)}\rangle\langle v^{(k)}|},
\end{split}
\end{equation}
the inner sum running over the $d$ elements $|v^{(k)}\rangle$ of $B^{(k)}$ and
$\oplus$ denoting addition modulo $d$. It is nonnegative on every
separable state \cite{Spengler}.

For comparison, \cite{BBCHM} displays the same operator as
\begin{gather*}
  \mathbf W^{\Gamma}(\mathcal M_m,s)
  =\frac{d+m-1}{d}\,\mathbf 1_d\otimes\mathbf 1_d-\mathbf B^{\Gamma}(\mathcal M_m,s),\\
  \mathbf B^{\Gamma}(\mathcal M_m,s)
  =\sum_{\ell=0}^{d-1}|\ell\rangle\langle\ell|\otimes|\ell+s\rangle\langle\ell+s|
   +\sum_{\alpha=1}^{m-1}\sum_{i=0}^{d-1}
    |i_{\alpha}\rangle\langle i_{\alpha}|\otimes
    |i_{\alpha}^{*}\rangle\langle i_{\alpha}^{*}| ,
\end{gather*}
where $\mathbf 1_d$ is the $d\times d$ identity, $\mathcal B_0$ is the canonical basis,
conjugation and transposition are taken with respect to $\mathcal B_0$, and $\ell+s$ is
read modulo $d$. Term by term this is
\eqref{eq:W}: the prefactor $\frac{d+m-1}{d}$, the shifted diagonal term and the
conjugation in the second tensor factor are theirs unchanged; our $k=2,\dots,m$ indexes
their $\alpha=1,\dots,m-1$ (the $m-1$ bases other than the computational one), our
$\overline{\,\cdot\,}$ is their ${}^{*}$, and our $\oplus$ is their $+$ modulo $d$.

\section{The objects}

\begin{equation}\label{eq:obj}
  \M_4=\{B^{(1)},\,B_{S_0},\,B_{S_{13}},\,B_{S_{19}}\},
  \qquad s=1,\qquad m=4=\tfrac d2,
\end{equation}
with $B^{(1)}$ the computational basis and, by \eqref{eq:idx},
\[
 S_0=\begin{pmatrix}0&0&0\\0&0&0\\0&0&0\end{pmatrix},\qquad
 S_{13}=\begin{pmatrix}1&0&1\\0&1&0\\1&0&0\end{pmatrix},\qquad
 S_{19}=\begin{pmatrix}1&1&0\\1&0&1\\0&1&0\end{pmatrix}.
\]

\begin{lemma}\label{lem:mub}
The four bases in \eqref{eq:obj} are orthonormal and pairwise mutually unbiased: all
$\binom42\cdot64=384$ cross overlaps satisfy $|\langle u\mid u'\rangle|^2=\frac18$.
\end{lemma}

Since $\sqrt8\,|v_S\rangle$ has entries in $\{\pm1,\pm\ii\}$, Lemma~\ref{lem:mub} is
an identity in $\mathbb Z[\ii]$: the $384$ overlaps of the rescaled vectors have
squared modulus exactly $8$, and the $64$ overlaps within a rescaled $B_S$ are
$8\delta_{vv'}$. No floating point occurs.

\begin{lemma}\label{lem:W}
$64\,\W(\M_4,1)$ is a symmetric matrix of integers of absolute value at most $64$; its
imaginary part vanishes identically. Its diagonal has $56$ entries equal to $64$ and
$8$ equal to $0$, so
\[
  \Tr \W(\M_4,1)=\tfrac{3584}{64}=56=d^2-d .
\]
\end{lemma}

Finally the state. Let $N=1049602=2\cdot524801$ and let $M\in\mathbb Z^{64\times64}$
be the symmetric matrix whose entries on and above the diagonal are those listed in
Table~\ref{tab:rho} and zero elsewhere, with $M_{ji}=M_{ij}$. Row $i$ of the table is
printed as \texttt{i | j:$M_{ij}$ j:$M_{ij}$ \dots} over the $j\ge i$ with
$M_{ij}\ne0$; the $288$ listed entries extend to $512$ nonzero entries of $M$. Put
\begin{equation}\label{eq:rho}
  \rho=\frac1N\,M .
\end{equation}
The indices are $i=8a+c$ with $a$ the index in the first tensor factor and $c$ that in
the second, so the partial transpose is
$(\rho^{\Gamma})_{8a+c,\,8b+e}=\rho_{8a+e,\,8b+c}$.

\begin{table}[htbp]
\caption{The $288$ integers on and above the diagonal of $M=N\rho$, $N=1049602$.}
\label{tab:rho}
\begin{center}
\begingroup\scriptsize
\begin{verbatim}
 0 | 0:13090 9:13229 18:13171 27:12881 36:12852 45:13179 54:13190 63:12939
 1 | 1:57650 8:13270 19:5884 26:4196 37:5429 44:4029 55:-6688 62:-5143
 2 | 2:12536 11:3191 16:3951 25:-5082 38:3182 47:12324 52:-5067 61:3955
 3 | 3:12118 10:4380 17:-6578 24:3461 39:-7078 46:5584 53:10793 60:4002
 4 | 4:12962 13:3941 22:4016 31:-4767 32:-4764 41:4021 50:3890 59:12707
 5 | 5:10924 12:4029 23:10749 30:4031 33:5426 40:-5093 51:5440 58:-5084
 6 | 6:13551 15:-6379 20:-5067 29:13165 34:3425 43:3589 48:4226 57:3866
 7 | 7:10281 14:-4675 21:5450 28:4002 35:10153 42:-4621 49:5587 56:4754
 8 | 8:3911 19:4196 26:1129 37:4029 44:1076 55:-5143 62:-1587
 9 | 9:13713 18:13534 27:13171 36:13179 45:13456 54:13398 63:13190
10 | 10:26484 17:4481 24:-1699 39:5584 46:-4404 53:4106 60:16959
11 | 11:12536 16:-5082 25:3951 38:12324 47:3182 52:3955 61:-5067
12 | 12:17079 23:4031 30:16880 33:-5093 40:1209 51:-5084 58:1239
13 | 13:13566 22:-5521 31:4016 32:4021 41:-5464 50:13385 59:3890
14 | 14:16021 21:4106 28:1086 35:-4621 42:15709 49:3658 56:891
15 | 15:13551 20:13165 29:-5067 34:3589 43:3425 48:3866 57:4226
16 | 16:14089 25:4956 38:-5067 47:3955 52:4930 61:13904
17 | 17:9972 24:3967 39:10793 46:4106 53:-6110 60:2457
18 | 18:13713 27:13229 36:13190 45:13398 54:13456 63:13179
19 | 19:57650 26:13270 37:-6688 44:-5143 55:5429 62:4029
20 | 20:13245 29:-3745 34:4226 43:3866 48:4411 57:4916
21 | 21:11631 28:-5748 35:5587 42:3658 49:9750 56:-5601
22 | 22:13566 31:3941 32:3890 41:13385 50:-5464 59:4021
23 | 23:10924 30:4029 33:5440 40:-5084 51:5426 58:-5093
24 | 24:15267 39:4002 46:16959 53:2457 60:-3588
25 | 25:14089 38:3955 47:-5067 52:13904 61:4930
26 | 26:3911 37:-5143 44:-1587 55:4029 62:1076
27 | 27:13090 36:12939 45:13190 54:13179 63:12852
28 | 28:30395 35:4754 42:891 49:-5601 56:15140
29 | 29:13245 34:3866 43:4226 48:4916 57:4411
30 | 30:17079 33:-5084 40:1239 51:-5093 58:1209
31 | 31:12962 32:12707 41:3890 50:4021 59:-4764
32 | 32:12962 41:3941 50:4016 59:-4767
33 | 33:10924 40:4029 51:10749 58:4031
34 | 34:13551 43:-6379 48:-5067 57:13165
35 | 35:10281 42:-4675 49:5450 56:4002
36 | 36:13090 45:13229 54:13171 63:12881
37 | 37:57650 44:13270 55:5884 62:4196
38 | 38:12536 47:3191 52:3951 61:-5082
39 | 39:12118 46:4380 53:-6578 60:3461
40 | 40:17079 51:4031 58:16880
41 | 41:13566 50:-5521 59:4016
42 | 42:16021 49:4106 56:1086
43 | 43:13551 48:13165 57:-5067
44 | 44:3911 55:4196 62:1129
45 | 45:13713 54:13534 63:13171
46 | 46:26484 53:4481 60:-1699
47 | 47:12536 52:-5082 61:3951
48 | 48:13245 57:-3745
49 | 49:11631 56:-5748
50 | 50:13566 59:3941
51 | 51:10924 58:4029
52 | 52:14089 61:4956
53 | 53:9972 60:3967
54 | 54:13713 63:13229
55 | 55:57650 62:13270
56 | 56:30395
57 | 57:13245
58 | 58:17079
59 | 59:12962
60 | 60:15267
61 | 61:14089
62 | 62:3911
63 | 63:13090
\end{verbatim}
\endgroup
\end{center}
\end{table}

\begin{lemma}\label{lem:rho}
$\rho$ defined by \eqref{eq:rho} is a state with positive partial transpose:
$\rho=\rho^{\top}$, $\Tr\rho=1$ (the $64$ diagonal entries of $M$ sum to $N$),
$\rho\succeq0$ and $\rho^{\Gamma}\succeq0$. Moreover
\[
  \Tr\bigl[\W(\M_4,1)\,\rho\bigr]
  =-\frac{14588}{524801}
  =-0.0277972031303\ldots<0 .
\]
\end{lemma}

\begin{proof}
Symmetry and the trace are immediate from Table~\ref{tab:rho}. For the two positivity
statements, the exact rational $LDL^{\!\top}$ factorisation without pivoting runs to
completion on $M$ and on $M^{\Gamma}$; all $64$ pivots are strictly positive in each
case, the smallest being
\begin{align*}
  \frac{63834948025622562}{457387529428466097047}
   &=1.395643\ldots\times10^{-4}\quad\text{for }\rho,\\
  \frac{329439466941518560}{2324566123790645652773}
   &=1.417208\ldots\times10^{-4}\quad\text{for }\rho^{\Gamma},
\end{align*}
after division by $N$. A Hermitian matrix admitting an $LDL^{\!\top}$ factorisation
with all pivots positive is positive definite. The trace is the integer pairing
$\sum_{i,j}(64\,\W)_{ij}M_{ji}$ divided by $64N$; the numerator is $-1867264$, and
$-1867264/(64\cdot1049602)=-14588/524801$. The largest denominator occurring in $\rho$
is $N$ itself, and every entry of $M$ has absolute value at most $57650$.
\end{proof}

\section{The result}

\begin{theorem}\label{thm:main}
$\W(\M_4,1)$ of \eqref{eq:obj} is non-decomposable. Consequently $m=4=\frac d2$
mutually unbiased bases of $\CC^8$ detect bound entanglement through \eqref{eq:W}, the
least number of mutually unbiased bases of $\CC^8$ for which \eqref{eq:W} has a
non-decomposable member is at most $4$, and the conjecture of \cite{BBCHM} quoted
in \S1, which asserts $\frac d2+1=5$ at $d=8$, is false at $r=3$.
\end{theorem}

\begin{proof}
Suppose $\W(\M_4,1)=P+Q^{\Gamma}$ with $P,Q\succeq0$. For any state $\sigma$ with
$\sigma^{\Gamma}\succeq0$,
\[
  \Tr\bigl[\W(\M_4,1)\sigma\bigr]
  =\Tr[P\sigma]+\Tr\bigl[Q^{\Gamma}\sigma\bigr]
  =\Tr[P\sigma]+\Tr\bigl[Q\,\sigma^{\Gamma}\bigr]\ \ge 0 ,
\]
because the trace of a product of two positive semidefinite operators is nonnegative
and $\Tr[XY^{\Gamma}]=\Tr[X^{\Gamma}Y]$. By Lemma~\ref{lem:rho} the state $\rho$ has
$\rho^{\Gamma}\succeq0$ and $\Tr[\W(\M_4,1)\rho]<0$, so no such $P,Q$ exist.

By \cite{Spengler} the operator $\W(\M_4,1)$ is nonnegative on separable
states, so $\Tr[\W(\M_4,1)\rho]<0$ forces $\rho$ to be entangled; since
$\rho^{\Gamma}\succeq0$, it is bound entangled. It is detected by the witness built
from the four MUBs of \eqref{eq:obj}, and Lemma~\ref{lem:mub} confirms that these are
four pairwise mutually unbiased bases. Hence four MUBs suffice at $d=8$, and the
minimum over the family \eqref{eq:W} is at most $4<5=\frac d2+1$.
\end{proof}

\section{Scope}

Only the cell $r=3$ is treated; dimensions $d=2^r$ with $r\ge4$ are untouched, and no
pattern is proposed in place of the refuted one. What is proved is ``at most $4$'', not
``exactly $4$'': whether $m=3$ or $m=2$ also suffice at $d=8$ is not addressed. Only the
particular $\M_4$ and the shift $s=1$ of \eqref{eq:obj} are certified, and nothing here
asserts that every set of four mutually unbiased bases of $\CC^8$, or every shift, gives
a non-decomposable witness. The number $-\frac{14588}{524801}$ is the value at the
$\rho$ printed above, not the minimum of $\Tr[\W(\M_4,1)\sigma]$ over PPT states
$\sigma$; only its sign is used in Theorem~\ref{thm:main}. That $\W(\M_m,s)$ is
nonnegative on separable states is the defining property established in
\cite{Spengler}; it is used, not reproved, here.

\section{Reproduction and verification}

Everything needed is printed above: the index conventions \eqref{eq:vS} and
\eqref{eq:idx}, the three symmetric matrices $S_0,S_{13},S_{19}$, the shift $s=1$,
the witness \eqref{eq:W}, and the $288$ integers of Table~\ref{tab:rho}. From these,
$64\,\W(\M_4,1)$ is assembled in $\mathbb Z[\ii]$ --- the rescaled basis vectors
$\sqrt8\,|v_S\rangle$ have entries in $\{\pm1,\pm\ii\}$, and each summand of
\eqref{eq:W} contributes $u_a\bar u_b\bar u_c u_e$ to the entry indexed by
$\bigl((a,c),(b,e)\bigr)$ --- and the four facts of Lemma~\ref{lem:rho} are three
integer computations and two rational $LDL^{\!\top}$ factorisations.

The calculation was re-run independently, in exact integer and rational arithmetic
throughout, from the objects as printed above and from nothing else. The program
rebuilds the four bases from \eqref{eq:vS}, checks orthonormality and the $384$ cross
overlaps of Lemma~\ref{lem:mub}, assembles $64\,\W(\M_4,1)$ and verifies
Lemma~\ref{lem:W}, parses Table~\ref{tab:rho}, and verifies Lemma~\ref{lem:rho}: the
trace, both $LDL^{\!\top}$ factorisations with their smallest pivots, and
$\Tr[\W(\M_4,1)\rho]$ by two independent routes. All of its checks passed. The program
and its transcript accompany this note; they are not needed in order to redo the
calculation, since every input they consume is printed above.

\begin{thebibliography}{9}

\bibitem{BBCHM}
J.~Bae, A.~Bera, D.~Chru\'sci\'nski, B.~C. Hiesmayr, and D.~McNulty,
\emph{How many mutually unbiased bases are needed to detect bound entangled states?},
J. Phys. A \textbf{55} (2022), no.~50, 505303.
\href{https://doi.org/10.1088/1751-8121/acaa16}{doi:10.1088/1751-8121/acaa16};
\href{https://arxiv.org/abs/2108.01109}{arXiv:2108.01109}.

\bibitem{Spengler}
C.~Spengler, M.~Huber, S.~Brierley, T.~Adaktylos, and B.~C. Hiesmayr,
\emph{Entanglement detection via mutually unbiased bases},
Phys. Rev. A \textbf{86} (2012), 022311.
\href{https://doi.org/10.1103/PhysRevA.86.022311}{doi:10.1103/PhysRevA.86.022311}.

\end{thebibliography}

\end{document}
